% ---------------------------------------------------------------------------
% Section 5 -- Evaluation
%
% Budget: 1,600-1,900 words + 5 tables (WRITING-GUIDE section 5).
% Structure: ~400 words of setup, then one 100-150 word reading per table.
% Bold run-in headings, no \subsection -- vertical space is the scarce resource.
%
% EVERY number here is copied verbatim from GROUND-TRUTH.md. Nothing is
% rounded, aggregated, or re-derived.
%
% Terminology: "resource" throughout, matching Sec.1, Sec.3 and Sec.4. The
% register rows the policy carries ARE the resources; GROUND-TRUTH.md calls
% them "assets", which is the artifact's word, not the paper's.
% ---------------------------------------------------------------------------
\section{Evaluation}
\label{sec:eval}

Four questions organise the evaluation. \textbf{RQ1}: can an organization's own
severity judgements be recovered from its policy prose alone? \textbf{RQ2}: do
the deterministic impact rules suffice, or does tool impact need a model?
\textbf{RQ3}: does the request-time signal add discrimination over the
design-time baseline? \textbf{RQ4}: does the calibration survive contact with
real traffic?

\textbf{Three corpora, one per question group.} RQ1 and RQ2 run on three vendor
MCP servers: a calendar server, GitHub, and Slack. Their tool catalogs are
genuine, captured from the live vendor servers through the protocol's own
\texttt{tools/list} call~\cite{mcp2026spec}, and are used exactly as returned
--- 13 tools for calendar, 26 for GitHub, 16 for Slack, 55 in total. RQ3 runs on
a different corpus: a session testbed spanning ten server profiles, whose
catalogs and whose sessions alike were written for this study, so RQ3 sits
entirely on the synthesized side of the accounting below and shares no catalog
with RQ1 or RQ2. RQ4 runs on a further server, a live \texttt{mcp-server-git}
driven over stdio, and is the one genuinely real execution trace in the section.

\textbf{What is synthesized.} Everything the captured catalogs do not supply.
The organizational policy document, the registers that name the resources behind
each of the three vendor servers (16, 20 and 20 resources, 56 in total), and the
ten RQ3 server profiles together with the invocation scenarios over them were
all written for this study. This is therefore an \emph{evaluation setting}, not
a benchmark, and we do not present it as one. It exists because graded
per-action severity ground truth for MCP invocations barely exists: we know of
no external benchmark that supplies it, which is also why RQ3 has no external
comparison point.

\textbf{Protocol for RQ1 and RQ2.} The scorer is given exactly two documents:
the captured tool catalog and the relevant section of the organizational policy.
The per-resource sensitivity table that ships with each server profile is
\emph{held out} and never enters the pipeline; it is the ground truth. That
table is the organization's own 1--5 rating of each register row, and the policy
document from which we derive our rating states no number anywhere. We report
mean absolute error, exact agreement, and agreement to within one tier. The
within-one column is reported because the sensitivity rating is not consumed on
its own: it is one of three 1--5 primitives whose product spans 1--125 before
banding (\S\ref{sec:framework}), so we evaluate each primitive against the
organization's judgement at the resolution the primitive is actually
authored at.

\textbf{RQ1: policy text recovers the organization's severities.}
Table~\ref{tab:sensitivity} gives the result per server. Across all 56
resources, 50 match the organization's own rating exactly, and \emph{none is off
by more than one tier} on any of the three servers. Mean absolute error is 0.125
on calendar and 0.10 on both GitHub and Slack. The uniform 100\% within-one
column is the load-bearing number: a derivation that reads only prose does not
merely correlate with the withheld table, it lands on it or immediately beside
it, everywhere. Exact agreement of 88--90\% is the weaker claim and the more
honest headline, because six resources do move.

\begin{table}[t]
  \centering
  \small
  \caption{RQ1 --- derived resource sensitivity against the organization's
    held-out per-resource table. The derivation reads the policy document only.}
  \label{tab:sensitivity}
  \begin{tabular}{lrrrr}
    \toprule
    Server & Resources & MAE & Exact & Within 1 \\
    \midrule
    calendar & 16 & 0.125 & 88\% & \textbf{100\%} \\
    github   & 20 & 0.10  & 90\% & \textbf{100\%} \\
    slack    & 20 & 0.10  & 90\% & \textbf{100\%} \\
    \bottomrule
  \end{tabular}
\end{table}

\textbf{Where it drifts.} Table~\ref{tab:misses} lists all six misses. Every one
is adjacent-tier and every one has a readable cause, which is the useful part.
Two are register-layout effects: \texttt{account-directory} inherits severity
from the Restricted-class account configuration it sits beside, and
\texttt{repository-catalog} is rated up because a full catalog enumeration maps
the estate. Two are classification-language effects: \texttt{infra-config} is
read as serious lasting harm rather than as control plane, and
\texttt{calendar-records} is read as a bare attribute listing. Two follow the
policy's own qualifiers: \texttt{incident-response} carries a
\texttt{self-sufficient} flag and credentials pasted mid-incident, and
\texttt{research-team} holds unpublished research read as competitively
damaging.
The misses run in both directions, the majority above the organization's own
rating rather than below --- the direction a security control should err in ---
though six cases establish a reading, not a calibrated bias.

\begin{table}[t]
  \centering
  \footnotesize
  \caption{RQ1 --- all six disagreements. Every miss is one tier, and every miss
    is traceable to a specific reading of the policy text.}
  \label{tab:misses}
  % \raggedright is applied per cell and rows end with \tabularnewline, so the
  % narrow text column does not need the array package (acmart does not load it).
  \begin{tabular}{@{}llrrp{2.6cm}@{}}
    \toprule
    Resource & Srv. & Ours & Org & Reading \tabularnewline
    \midrule
    \texttt{account-directory} & cal & 5 & 4 & \raggedright sits beside the Restricted-class account configuration \tabularnewline
    \texttt{calendar-records}  & cal & 2 & 3 & \raggedright read as a bare attribute listing \tabularnewline
    \texttt{infra-config}      & gh  & 4 & 5 & \raggedright serious lasting harm, not control plane \tabularnewline
    \texttt{repository-catalog}& gh  & 3 & 2 & \raggedright full enumeration maps the estate \tabularnewline
    \texttt{incident-response} & sl  & 5 & 4 & \raggedright \texttt{self-sufficient} flag, credentials mid-incident \tabularnewline
    \texttt{research-team}     & sl  & 4 & 3 & \raggedright unpublished research, competitively damaging \tabularnewline
    \bottomrule
  \end{tabular}
\end{table}

\textbf{RQ2: the impact ladder needed no model.} Tool impact is the one stage
where a model is optional by construction --- the rules answer, and the model is
consulted only when the rules abstain (\S\ref{sec:framework}).
Table~\ref{tab:impact} reports what happened. The deterministic ladder decided
all 55 tools across the three servers and the LLM fallback was invoked zero
times. Compared against an earlier arm in which tool impact was scored by the
LLM, the ladder agrees on 51 of 55 tools (93\%), and all four disagreements are
one tier: \texttt{get\_pull\_request\_files} (3 vs.\ 2), \texttt{search\_users}
(3 vs.\ 2), \texttt{conversations\_leave} (4 vs.\ 5) and \texttt{usergroups\_me}
(5 vs.\ 4). The consequence is that the tool axis can be decided without a model
call, at 93\% agreement with the LLM arm it replaces and with every
disagreement confined to one tier.

\begin{table}[t]
  \centering
  \small
  \caption{RQ2 --- tool impact. The rules answered every tool; the LLM fallback
    never fired. Agreement is against an earlier arm in which the LLM scored
    tool impact.}
  \label{tab:impact}
  \begin{tabular}{lrrrr}
    \toprule
    Server & Tools & Ladder & LLM & Agrees \\
    \midrule
    calendar & 13 & 13 & 0 & 13/13 \\
    github   & 26 & 26 & 0 & 24/26 \\
    slack    & 16 & 16 & 0 & 14/16 \\
    \bottomrule
  \end{tabular}
\end{table}

\textbf{An honesty constraint on RQ2.} The fallback did not fire, and the reason
matters more than the result. Every one of the 55 tools matched a tier verb, so
every rule confidence was 0.8. The 0.95 value was unreachable, because none of
the three vendor catalogs declares MCP tool annotations; and no tool fell to the
0.35 abstention value that would have handed it to the model
(\S\ref{sec:handoff}). On this corpus the rules-to-model hand-off is therefore
covered by unit tests only. We claim that the ladder was sufficient here; we do
not claim the hybrid was exercised.

\textbf{RQ3: the request-time signal adds discrimination.} The scenario
generator partitions sessions three ways. Adversarial sessions are the target
class; benign sessions are the population that must not be disturbed; and
\emph{misuse} sessions are accidental-insider behaviour, constructed to sit
between the two, which gives the measurement an interior class rather than only
a separable pair. The ten server profiles contribute 86--498 benign and 4--20
adversarial sessions each, 88 adversarial sessions in total.
Table~\ref{tab:dynamic} reports the two operating thresholds an organization
would realistically choose, high and critical. Each session is scored twice ---
static only, and static combined with the request-time signal --- so the
difference is exactly what the dynamic layer contributes; a session takes the
band of its riskiest call.

The critical row is the result. Holding benign fall-out fixed at 9\%,
adversarial recall rises from 56\% to 81\%, and 22 of the 88 adversarial
sessions are caught \emph{only} by the dynamic signal. That is 25 points of
recall bought at no cost in false alarms, which is the shape an operator needs:
the static scorer alone reaches comparable recall only by dropping to the
permissive threshold, where it recovers 91\% but flags 49\% of benign sessions.
At that permissive threshold the dynamic layer still helps and is merely
saturated, 91\% to 100\% recall with 8 of 88 caught only dynamically.

\textbf{The interior class does not move at the strict threshold.} We built
misuse in as the hard case, so its row is the one to read honestly. At the
critical threshold the dynamic layer buys nothing on it --- 24\% flagged by the
static scorer, 24\% by the combined one --- in the same sweep where adversarial
recall gains 25 points. Misuse improves only at the permissive threshold, 63\%
to 71\%. On this testbed the request-time signal separates deliberate
adversarial behaviour and leaves accidental-insider behaviour where the static
scorer left it. We did not run the ablation that would attribute this to a
particular signal, so we report it as a bound on the claim rather than explain
it away: the layer buys adversarial recall, not insider recall.

Mean session band-rank (low\,=\,1 to critical\,=\,4) moves from static to
combined as follows: adversarial 3.47 to 3.81, misuse 2.84 to 2.92, benign 2.41
to 2.41. Benign sessions are unchanged on both measures, band-rank and
fall-out, which is the property an operator needs before turning the layer on.
Two design constraints bound how far the layer can move anything at all: the
band signals can only raise risk above the static floor, and the likelihood
multiplier is bounded below at 0.1, so no amount of apparent normality discounts
a critical cell to nothing.

\begin{table}[t]
  \centering
  \small
  \caption{RQ3 --- session-level discrimination, static versus static combined
    with the request-time signal, over ten synthesized server profiles (86--498
    benign and 4--20 adversarial sessions each).}
  \label{tab:dynamic}
  \begin{tabular}{lrrr}
    \toprule
    Scorer & Recall (adv.) & Flagged (misuse) & Fall-out (benign) \\
    \midrule
    \multicolumn{4}{@{}l}{\emph{Flag threshold = high}} \\
    static   & 91\%           & 63\% & 49\% \\
    combined & \textbf{100\%} & 71\% & 49\% \\
    \midrule
    \multicolumn{4}{@{}l}{\emph{Flag threshold = critical}} \\
    static   & 56\%          & 24\% & 9\% \\
    combined & \textbf{81\%} & 24\% & \textbf{9\%} \\
    \bottomrule
  \end{tabular}
\end{table}

\textbf{RQ4: the calibration transfers to real traffic.} Everything in RQ3 is
synthetic, so we repeated the calibration check on genuinely real MCP traffic.
We drove a live \texttt{mcp-server-git} over stdio against the implementation's
own repository (223 real commits) and captured 400 real tool calls, of which 399
executed without error, exercising 12 git tools and 106 distinct call
signatures. Arguments come from the repository's actual history: real 40-char
SHAs, real paths, real branch names. All 400 captured calls are benign by
construction, being script-driven replay against that history; the only
non-benign items in RQ4 are three calls we injected. We fitted on the first 70\%
of the calls chronologically and scored the held-out 30\%. The false-positive
rate on benign traffic was 0.8\%, against a design target of 1--2\%, and the
mean held-out likelihood was 0.107 --- essentially everything sits on the
multiplier's 0.1 floor, which is the intended behaviour for traffic that looks
like its own history. The three injected calls were a 200-char commit message, a
bare \texttt{git\_reset}, and a fabricated \texttt{exfiltrate\_all} carrying a
payload; all three scored high, against held-out benign traffic sitting on that
floor. The q99 quantile-anchored ramp was tuned on the synthetic calendar,
GitHub and Slack corpus and transferred to real git traffic with no adjustment.

\textbf{Limitations of the measurement.} Blast radius is the noisy stage.
Table~\ref{tab:blast} compares the (tool, resource) cells produced with and
without the policy context: 19--22\% of cells move, with roughly half the
movement being the relevance gate flipping a cell into or out of
not-applicable. An earlier arm measured 23--35 cells per server of pure
run-to-run variance with the prompt held fixed. Calendar's 40 changed cells sit
just above that band, but GitHub's 100 and Slack's 72 are well above it, so on
the two larger servers the movement cannot be attributed to run-to-run variance
alone; we nonetheless do not read individual cell changes as signal. The
movement is not a uniform discount either: summed scores rise on calendar (2280
to 2697) and Slack (3779 to 4096) and fall on GitHub (6340 to 5328). Three
further limits bound the claims. The RQ3 labels are synthetic by construction
--- the benign, misuse and adversarial partition is a property of the generator,
not an independent annotation --- so RQ3 measures separation under a known
generative model, and RQ4 only checks that the calibration holds elsewhere. The
optional LLM judge in the request-time layer was wired but not run in these
experiments. And the static result rests on one policy document over three
servers, so cross-organization generalisation is untested; beyond catalog
capture, the vendor servers themselves were not driven live.

\begin{table}[t]
  \centering
  \small
  \caption{Blast radius is the noisy stage: (tool, resource) cells that changed
    between arms under the same rubric and the same resource identifiers.}
  \label{tab:blast}
  \begin{tabular}{lrrrrrr}
    \toprule
    Server & Cells & Diff. & Raised & Lower. & $\to$N/A & N/A$\to$ \\
    \midrule
    calendar & 208 & 40  & 11 & 3  & 15 & 11 \\
    github   & 520 & 100 & 26 & 21 & 32 & 21 \\
    slack    & 320 & 72  & 30 & 18 & 16 & 8  \\
    \bottomrule
  \end{tabular}
\end{table}
